@@M0@@ be the class of such a sheaf. We then have:
@@M1@@
Since @@M2@@, the definition of @@M3@@ gives:
@@M4@@

Furthermore, @@M5@@ has augmentation @@M6@@; the Chern class occurring in formula 4.4 for @@M7@@ is therefore @@M8@@. Taking into account the definition of @@M9@@, and 3.5, the relation to be proved is written:
@@M10@@
and is therefore verified.

\sgaStableTarget{sga6.e08.rem.4-4-1}\textbf{Remark 4.4.1.} It is also possible to deduce 4.4 from a corresponding purely algebraic statement, cf. @@M11@@ I 1.3.

\sgaStableTarget{sga6.e08.para.4-5}\textbf{4.5.} We shall now show that, to prove the Riemann--Roch theorem in the case considered in this section, one may reduce to the case where @@M12@@ is a divisor of @@M13@@, and where the element to which the formula applies is the inverse image of an element of @@M14@@.

By 4.2, the formula to be proved:
@@M15@@
is equivalent to the formula:
@@M16@@
Taking into account the functoriality of @@M17@@, and 4.3, it may also be written:
@@M18@@
or, by VII 3.4:
@@M19@@
